\chapter{Topologie na rodzinach podzbiorów}
\label{ch:topog-hiperprzestrzenie}

\noindent\textbf{Cel rozdziału.}
W rozdziale~\ref{ch:topog-przestrzenie-odwzorowan} rozważaliśmy
topologie na zbiorach odwzorowań. Teraz punktami nowej przestrzeni
będą same podzbiory ustalonej przestrzeni \(X\).
Pozwala to badać rodziny zbiorów zwartych, zmieniające się
figury geometryczne i granice zbiorów domkniętych.
Rozróżnimy topologię Vietorisa na zbiorach zwartych
oraz topologię Fella na zbiorach domkniętych.
W przypadku metrycznym wykażemy związek z metryką Hausdorffa,
udowodnimy zupełność odpowiedniej przestrzeni zbiorów zwartych
i wskażemy, kiedy dwie topologie opisują różne rodzaje zbieżności.

\section{Hiperprzestrzeń zbiorów zwartych i topologia Vietorisa}
\label{sec:topog21-vietoris}

\begin{definicja}[Hiperprzestrzeń zwartych podzbiorów]
\label{def:topog21-kx}
Dla przestrzeni topologicznej \(X\) definiujemy
\(\mathcal K(X)=\{K\subseteq X:K\neq\varnothing,\ K
\text{ jest zwarty}\}\).
Dla zbioru otwartego \(U\subseteq X\) wprowadzamy
\[
 U^-=\{K\in\mathcal K(X):K\cap U\neq\varnothing\},
 \qquad
 U^+=\{K\in\mathcal K(X):K\subseteq U\}.
\]
\emph{Topologią Vietorisa} na \(\mathcal K(X)\)
nazywamy topologię generowaną przez wszystkie zbiory
\(U^-\) i \(U^+\), gdzie \(U\) przebiega zbiory otwarte w \(X\).
Zbiory pierwszego rodzaju określają warunki
\emph{trafienia} zbioru otwartego, a drugiego
warunki \emph{zawierania się} w zbiorze otwartym.
\end{definicja}

\begin{uwaga}[Rola zbioru pustego]
\label{rem:topog21-empty}
Zbiór pusty został wyłączony z \(\mathcal K(X)\).
Dzięki temu zarówno odległość Hausdorffa,
jak i warunek trafienia mają zwykłą postać,
bez wartości nieskończonych i dodatkowych wyjątków.
W topologii Fella będziemy natomiast dopuszczać zbiór pusty.
Jeśli \(X=\varnothing\), to \(\mathcal K(X)=\varnothing\).
\end{uwaga}

\begin{stwierdzenie}[Baza topologii Vietorisa]
\label{prop:topog21-vietoris-basis}
Niech \(X\) będzie przestrzenią topologiczną.
Dla każdego \(m\ge1\) i otwartych \(U_1,\ldots,U_m\subseteq X\)
połóżmy
\[
 \langle U_1,\ldots,U_m\rangle
 =\{K\in\mathcal K(X):
 K\subseteq\textstyle\bigcup_{i=1}^{m}U_i,\
 K\cap U_i\neq\varnothing\text{ dla }1\le i\le m\}.
\]
Rodzina wszystkich takich zbiorów jest bazą
topologii Vietorisa.
\end{stwierdzenie}
\begin{proof}
Każdy rozważany zbiór ma postać
\((\bigcup_iU_i)^+\cap\bigcap_iU_i^-\),
a więc jest otwarty.
Aby otrzymać bazę, wystarczy dla dowolnego
\(K\in\mathcal K(X)\) oraz dowolnego
skończonego przecięcia zbiorów podbazowych,
które zawiera \(K\), znaleźć zbiór przedstawionej
postaci zawierający \(K\) i zawarty w tym przecięciu.

Niech warunki górne mają postać
\(K\subseteq V_1,\ldots,K\subseteq V_r\),
a dolne \(K\cap W_j\neq\varnothing\) dla \(1\le j\le s\).
Przyjmijmy \(V=\bigcap_{i=1}^{r}V_i\),
a dla \(r=0\) przyjmijmy \(V=X\).
Wówczas
\(\langle V,V\cap W_1,\ldots,V\cap W_s\rangle\)
zawiera \(K\): warunki trafienia \(V\cap W_j\)
wynikają z \(K\subseteq V\) oraz trafienia \(W_j\).
Każdy element tego zbioru jest zawarty w \(V\)
i trafia wszystkie \(W_j\).
Dotyczy to także \(s=0\), gdy bazowym zbiorem
jest po prostu \(\langle V\rangle=V^+\).
\end{proof}

\begin{stwierdzenie}[Jednopunktowe zbiory jako podprzestrzeń]
\label{prop:topog21-singletons}
Niech \(X\) będzie przestrzenią topologiczną.
Odwzorowanie
\(j\colon X\to\mathcal K(X)\), \(j(x)=\{x\}\),
jest zanurzeniem topologicznym.
Jeżeli ponadto \(X\) jest przestrzenią Hausdorffa,
to \(j(X)\) jest domknięty w \(\mathcal K(X)\).
\end{stwierdzenie}
\begin{proof}
Dla każdego otwartego \(U\subseteq X\) mamy
\(j^{-1}(U^-)=j^{-1}(U^+)=U\).
Stąd \(j\) jest ciągłe.
Jest różnowartościowe, a odwzorowanie odwrotne
na jego obraz jest ciągłe, ponieważ
\(j(U)=j(X)\cap U^-\) jest otwarte w \(j(X)\).
Jeżeli \(X\) jest Hausdorffa, a \(K\in\mathcal K(X)\)
ma przynajmniej dwa punkty \(x\neq y\), to istnieją
rozłączne zbiory otwarte \(U\ni x\), \(V\ni y\).
Zbiór \(U^-\cap V^-\) jest otwartym otoczeniem \(K\),
niezawierającym żadnego singletonu.
Zatem dopełnienie \(j(X)\) jest otwarte.
\end{proof}

\begin{twierdzenie}[Oddzielanie zbiorów zwartych]
\label{thm:topog21-vietoris-separation}
Jeżeli \(X\) jest przestrzenią Hausdorffa,
to \(\mathcal K(X)\) z topologią Vietorisa
również jest przestrzenią Hausdorffa.
\end{twierdzenie}
\begin{proof}
Niech \(A,B\in\mathcal K(X)\) i \(A\neq B\).
Po ewentualnej zamianie ich ról wybierzmy \(x\in A\setminus B\).
Dla każdego \(b\in B\) istnieją rozłączne zbiory otwarte
\(U_b\ni x\) i \(V_b\ni b\).
Zwartość \(B\) pozwala wybrać \(b_1,\ldots,b_m\)
tak, że \(B\subseteq\bigcup_{i=1}^m V_{b_i}\).
Wówczas \(U=\bigcap_{i=1}^m U_{b_i}\) jest otwartym
otoczeniem \(x\), a \(V=\bigcup_{i=1}^mV_{b_i}\)
jest otwartym otoczeniem \(B\), przy czym
\(U\cap V=\varnothing\).
Zbiory \(U^-\) i \(V^+\) są rozłącznymi otwartymi
otoczeniami odpowiednio \(A\) i \(B\)
w \(\mathcal K(X)\).
\end{proof}

\begin{stwierdzenie}[Obrazy zwartych zbiorów]
\label{prop:topog21-induced}
Niech \(X,Y\) będą przestrzeniami topologicznymi,
a \(f\colon X\to Y\) odwzorowaniem ciągłym.
Wówczas odwzorowanie
\[
 \mathcal K(f)\colon\mathcal K(X)\to\mathcal K(Y),
 \qquad K\longmapsto f(K),
\]
jest dobrze określone i ciągłe dla topologii Vietorisa.
\end{stwierdzenie}
\begin{proof}
Obraz niepustego zbioru zwartego przez odwzorowanie
ciągłe jest niepusty i zwarty.
Dla otwartego \(U\subseteq Y\) zachodzą równości
\[
 \mathcal K(f)^{-1}(U^-)=\bigl(f^{-1}(U)\bigr)^-,
 \qquad
 \mathcal K(f)^{-1}(U^+)=\bigl(f^{-1}(U)\bigr)^+.
\]
Oba przeciwobrazy są otwarte, ponieważ \(f^{-1}(U)\)
jest otwarty w \(X\). Wystarczy sprawdzić
ciągłość na podbazie.
\end{proof}

\section{Metryka Hausdorffa i jej topologia}
\label{sec:topog21-hausdorff}

\begin{definicja}[Odległość Hausdorffa]
\label{def:topog21-hausdorff}
Niech \((X,d)\) będzie przestrzenią metryczną.
Dla \(x\in X\) i niepustego \(A\subseteq X\) piszemy
\(d(x,A)=\inf_{a\in A}d(x,a)\).
Dla \(A,B\in\mathcal K(X)\) definiujemy
\[
 d_H(A,B)=
 \max\left\{\sup_{a\in A}d(a,B),
             \sup_{b\in B}d(b,A)\right\}.
\]
Obie liczby pod supremami są skończone,
gdyż \(A\) i \(B\) są niepuste i zwarte,
a więc ograniczone.
\end{definicja}

\begin{twierdzenie}[Metryczność odległości Hausdorffa]
\label{thm:topog21-hausdorff-metric}
Dla każdej przestrzeni metrycznej \((X,d)\)
funkcja \(d_H\) jest metryką na \(\mathcal K(X)\).
\end{twierdzenie}
\begin{proof}
Nieujemność i symetria są bezpośrednie.
Jeżeli \(d_H(A,B)=0\), to dla każdego \(a\in A\)
mamy \(d(a,B)=0\). Ponieważ \(B\) jest zwarty,
jest domknięty w \(X\), stąd \(a\in B\).
Otrzymujemy \(A\subseteq B\), a analogicznie
\(B\subseteq A\), więc \(A=B\).

Niech \(A,B,C\in\mathcal K(X)\).
Dla \(a\in A\) i \(b\in B\) nierówność trójkąta
daje \(d(a,C)\le d(a,b)+d(b,C)\).
Biorąc infimum po \(b\in B\), otrzymujemy
\(d(a,C)\le d(a,B)+\sup_{b\in B}d(b,C)\).
Po wzięciu supremum po \(a\in A\) wynika
\[
 \sup_{a\in A}d(a,C)
 \le d_H(A,B)+d_H(B,C).
\]
Identycznie oszacowujemy \(\sup_{c\in C}d(c,A)\).
Maksimum obu liczb spełnia więc
\(d_H(A,C)\le d_H(A,B)+d_H(B,C)\).
\end{proof}

\begin{przyklad}[Dwa skończone podzbiory prostej]
\label{ex:topog21-points}
Dla \(A=\{0,2\}\) i
\(B=\{\frac15,\frac{17}{10}\}\)
w metryce euklidesowej mamy
\(\sup_{a\in A}d(a,B)=\frac3{10}\) oraz
\(\sup_{b\in B}d(b,A)=\frac3{10}\);
zatem \(d_H(A,B)=\frac3{10}\).
Nie jest to odległość między środkami
ani suma odległości dobranych punktów.
\begin{figure}[htbp]
\centering
\begin{tikzpicture}[x=3.05cm,y=1cm]
 \draw[-{Latex[length=2mm]},thin] (-.13,0)--(2.25,0);
 \foreach \t/\lab in {0/0,.2/{0{,}2},1.7/{1{,}7},2/2}
   {\draw (\t,.04)--(\t,-.04);
    \node[below,font=\small] at (\t,-.08) {$\lab$};}
 \fill[blue!70!black] (0,.52) circle (2.6pt);
 \fill[blue!70!black] (2,.52) circle (2.6pt);
 \fill[orange!85!black] (.2,.28) circle (2.6pt);
 \fill[orange!85!black] (1.7,.28) circle (2.6pt);
 \node[left,font=\small,blue!70!black] at (-.12,.67) {$A$};
 \node[left,font=\small,orange!85!black] at (-.12,.13) {$B$};
\end{tikzpicture}
\caption{Odległość Hausdorffa mierzy najdalszą odległość
punktu jednego zbioru od drugiego zbioru.}
\label{fig:topog21-points}
\end{figure}
\end{przyklad}

\begin{twierdzenie}[Vietoris i metryka Hausdorffa]
\label{thm:topog21-vietoris-hausdorff}
Niech \((X,d)\) będzie przestrzenią metryczną.
Topologia Vietorisa na \(\mathcal K(X)\) jest równa
topologii metryki \(d_H\).
\end{twierdzenie}
\begin{proof}
Najpierw sprawdzimy, że wszystkie zbiory podbazowe
Vietorisa są otwarte w metryce \(d_H\).
Jeżeli \(K\cap U\neq\varnothing\) i \(U\subseteq X\)
jest otwarty, wybieramy \(x\in K\cap U\)
oraz \(r>0\) z \(B_d(x,r)\subseteq U\).
Gdy \(d_H(K,L)<r/2\), istnieje \(y\in L\)
z \(d(x,y)<r/2\), więc \(L\cap U\neq\varnothing\).
Zatem \(U^-\) jest otwarty.

Jeżeli \(K\subseteq U\), to dla \(U=X\) teza jest
oczywista. W przeciwnym razie funkcja
\(x\mapsto d(x,X\setminus U)\) jest ciągła
i dodatnia na zwartym \(K\). Osiąga więc
dodatnie minimum \(r\).
Jeżeli \(d_H(K,L)<r/2\), to każdy \(y\in L\)
ma odległość mniejszą niż \(r/2\) od pewnego
\(x\in K\), zatem \(y\notin X\setminus U\).
Daje to \(L\subseteq U\) i otwartość \(U^+\).

Odwrotnie, niech \(K\in\mathcal K(X)\) i \(r>0\).
Zwartość \(K\) pozwala wybrać \(x_1,\ldots,x_m\in K\)
tak, że \(K\subseteq\bigcup_i B_d(x_i,r/3)\).
Niech \(V=\bigcup_{x\in K}B_d(x,r/2)\).
Rozważmy otoczenie Vietorisa
\[
 \mathcal O=V^+\cap
 \bigcap_{i=1}^{m} B_d(x_i,r/3)^- .
\]
Zbiór \(K\) należy do \(\mathcal O\).
Dla \(L\in\mathcal O\), każdego \(y\in L\)
mamy \(d(y,K)<r/2\), ponieważ \(L\subseteq V\).
Dla dowolnego \(x\in K\) wybieramy \(i\)
z \(d(x,x_i)<r/3\), a następnie
\(y_i\in L\cap B_d(x_i,r/3)\).
Wtedy \(d(x,L)<2r/3\).
Wobec tego \(d_H(K,L)<r\),
czyli \(\mathcal O\subseteq B_{d_H}(K,r)\).
Otrzymujemy drugie zawieranie topologii.
\end{proof}

\begin{wniosek}[Zbieżność w metryce Hausdorffa]
\label{cor:topog21-hausdorff-convergence}
Niech \((X,d)\) będzie przestrzenią metryczną
i \(K_n,K\in\mathcal K(X)\).
Wówczas \(K_n\to K\) w topologii Vietorisa
dokładnie wtedy, gdy dla każdego \(\varepsilon>0\)
od pewnego miejsca zachodzą oba zawierania
\[
 K_n\subseteq\{x:d(x,K)<\varepsilon\},
 \qquad
 K\subseteq\{x:d(x,K_n)<\varepsilon\}.
\]
\end{wniosek}
\begin{proof}
Oba zawierania są równoważne nierówności
\(d_H(K_n,K)\le\varepsilon\) z odpowiednią
zmianą \(\varepsilon\) na mniejszą liczbę dodatnią.
Zastosujmy Twierdzenie~\ref{thm:topog21-vietoris-hausdorff}
i definicję granicy ciągu w przestrzeni metrycznej.
\end{proof}

\section{Zupełność i zwartość hiperprzestrzeni}
\label{sec:topog21-complete}

\begin{twierdzenie}[Zupełność hiperprzestrzeni]
\label{thm:topog21-complete}
Jeżeli \((X,d)\) jest zupełną przestrzenią metryczną,
to \((\mathcal K(X),d_H)\) jest zupełną przestrzenią
metryczną.
\end{twierdzenie}
\begin{proof}
Niech \((K_n)\) będzie ciągiem Cauchy'ego
w \((\mathcal K(X),d_H)\).
Wybierzmy rosnący ciąg indeksów \((n_j)\)
tak, że dla każdego \(j\ge1\)
\[
 d_H(K_{n_j},K_{n_{j+1}})<2^{-j}.
\]
Taki wybór jest możliwy z warunku Cauchy'ego.
Oznaczmy \(A_j=K_{n_j}\)
i \(f_j(x)=d(x,A_j)\) dla \(x\in X\).
Nierówność
\(|d(x,A)-d(x,B)|\le d_H(A,B)\)
pokazuje, że \((f_j)\) jest jednostajnie Cauchy'ego
na \(X\). Zatem istnieje funkcja \(f\colon X\to[0,\infty)\),
do której \(f_j\) zbiega jednostajnie.
Każde \(f_j\) jest \(1\)-lipschitzowskie,
więc \(f\) jest ciągła.
Niech \(K=\{x\in X:f(x)=0\}\).
Zbiór \(K\) jest domknięty.

Pokażemy, że \(K\neq\varnothing\).
Wybieramy \(x_1\in A_1\);
rekurencyjnie można wybrać \(x_{j+1}\in A_{j+1}\)
tak, aby \(d(x_j,x_{j+1})<2^{-j}\).
Wówczas \((x_j)\) jest ciągiem Cauchy'ego,
więc z zupełności \(X\) zbiega do \(x\in X\).
Ponieważ \(d(x,A_j)\le d(x,x_j)\to0\),
a \(f_j(x)\to f(x)\), mamy \(f(x)=0\),
czyli \(x\in K\).

Dla dowolnego \(j\) i \(z\in A_j\)
można analogicznie rozpocząć łańcuch
w punkcie \(z\), otrzymując punkt
\(x_z\in K\) z
\(d(z,x_z)\le\sum_{k=j}^{\infty}2^{-k}=2^{1-j}\).
Zatem \(\sup_{z\in A_j}d(z,K)\le 2^{1-j}\).
Z drugiej strony dla \(x\in K\) mamy
\(d(x,A_j)=|f_j(x)-f(x)|
\le\sup_{y\in X}|f_j(y)-f(y)|\to0\).
Po wykazaniu zwartości \(K\)
oszacowania te pozwolą zastosować
definicję \(d_H\) i otrzymać \(d_H(A_j,K)\to0\).
Udowodnimy więc zwartość \(K\)
bez zakładania jej z góry.

Niech \(\varepsilon>0\).
Wybierzmy \(j\) tak, że
\(\sup_{x\in K}d(x,A_j)<\varepsilon/3\).
Zwarty \(A_j\) ma skończoną
\(\varepsilon/3\)-sieć \(z_1,\ldots,z_m\in A_j\).
Każdy punkt \(x\in K\) leży
w odległości mniejszej niż \(2\varepsilon/3\)
od któregoś \(z_i\).
Zatem \(K\) jest całkowicie ograniczony.
Jako domknięty podzbiór zupełnej przestrzeni
\(X\) jest również zupełny, a więc zwarty.
Dowiedzione oszacowania dają
\(d_H(A_j,K)\to0\).
Cały ciąg Cauchy'ego \((K_n)\)
zbiega wtedy do \(K\), ponieważ ma podciąg
zbieżny do \(K\).
\end{proof}

\begin{twierdzenie}[Zwartość dla zwartej przestrzeni metrycznej]
\label{thm:topog21-compact-metric}
Jeżeli \((X,d)\) jest zwartą przestrzenią metryczną,
to \((\mathcal K(X),d_H)\) jest zwarte.
\end{twierdzenie}
\begin{proof}
Przestrzeń \(X\) jest zupełna,
więc Twierdzenie~\ref{thm:topog21-complete}
zapewnia zupełność \(\mathcal K(X)\).
Wykażemy całkowitą ograniczoność tej hiperprzestrzeni.

Niech \(\varepsilon>0\).
Zwartość \(X\) daje skończoną
\(\varepsilon/3\)-sieć \(S\subseteq X\).
Dla każdego \(K\in\mathcal K(X)\) połóżmy
\(S_K=\{s\in S:d(s,K)<\varepsilon/3\}\).
Zbiór \(S_K\) jest niepusty: każdy \(x\in K\)
ma odległość mniejszą niż \(\varepsilon/3\)
od pewnego punktu \(S\).
Dla \(x\in K\) możemy wybrać \(s\in S_K\)
z \(d(x,s)<\varepsilon/3\); dla \(s\in S_K\)
z definicji \(d(s,K)<\varepsilon/3\).
Zatem \(d_H(K,S_K)\le\varepsilon/3<\varepsilon\).
Ponieważ niepustych podzbiorów skończonego \(S\)
jest skończenie wiele, stanowią one skończoną
\(\varepsilon\)-sieć w \(\mathcal K(X)\).
Zupełność i całkowita ograniczoność
metrycznej \(\mathcal K(X)\) dają jej zwartość.
\end{proof}

\begin{twierdzenie}[Zwartość Vietorisa bez metryki]
\label{thm:topog21-vietoris-compact}
Jeżeli \(X\) jest zwartą przestrzenią Hausdorffa,
to \(\mathcal K(X)\) z topologią Vietorisa
jest zwartą przestrzenią Hausdorffa.
\end{twierdzenie}
\begin{proof}
Własność Hausdorffa wynika z
Twierdzenia~\ref{thm:topog21-vietoris-separation}.
Aby dowieść zwartości, zastosujemy
charakterystykę przez ultrafiltry z
Twierdzenia~\ref{thm:topog-zwartosc-ultrafiltry}.
Niech \(\mathcal U\) będzie ultrafiltrem
na \(\mathcal K(X)\).
Określmy
\[
 K=\{x\in X:\text{ dla każdego otwartego }
  V\ni x\text{ zachodzi }V^-\in\mathcal U\}.
\]
Jeśli \(x\notin K\), istnieje otwarte \(V\ni x\)
z \(V^-\notin\mathcal U\). Każdy punkt \(V\)
również nie należy do \(K\).
Zatem \(K\) jest domknięty w \(X\).

Gdyby \(K\) był pusty, moglibyśmy dla każdego
\(x\in X\) wybrać otwarty \(V_x\ni x\)
z \(V_x^-\notin\mathcal U\).
Zwartość \(X\) pozwala wybrać skończone
podpokrycie \(V_{x_1},\ldots,V_{x_m}\).
Każdy niepusty zbiór zwarty trafia
przynajmniej jedno \(V_{x_i}\), czyli
\(\mathcal K(X)=\bigcup_i V_{x_i}^-\).
Ultrafiltr musiałby zawierać
co najmniej jeden \(V_{x_i}^-\),
co daje sprzeczność.
Stąd \(K\neq\varnothing\), a jako zbiór
domknięty w zwartej \(X\) jest zwarty.

Sprawdzimy, że \(\mathcal U\to K\) w topologii Vietorisa.
Jeśli \(K\cap V\neq\varnothing\) dla otwartego \(V\),
to z definicji \(K\) mamy \(V^-\in\mathcal U\).
Jeśli \(K\subseteq V\), zbiór \(C=X\setminus V\)
jest zwarty. Dla każdego \(x\in C\) wybieramy
otwarty \(W_x\ni x\) z \(W_x^-\notin\mathcal U\).
Z pokrycia \(C\) wybieramy skończenie wiele
\(W_{x_1},\ldots,W_{x_m}\).
Dopełnienie \(V^+\) w \(\mathcal K(X)\)
jest zawarte w \(\bigcup_iW_{x_i}^-\).
Nie może należeć do ultrafiltru, gdyż wtedy
jedno z \(W_{x_i}^-\) również do niego należałoby.
Zatem \(V^+\in\mathcal U\).
Ultrafiltr zawiera każde podbazowe otoczenie
\(K\), a więc każde otoczenie \(K\).
Jest zbieżny do \(K\).
\end{proof}

\section{Topologia Fella na zbiorach domkniętych}
\label{sec:topog21-fell}

\begin{definicja}[Hiperprzestrzeń Fella]
\label{def:topog21-fell}
Niech \(X\) będzie przestrzenią topologiczną.
Przez \(\mathcal F(X)\) oznaczamy rodzinę
\emph{wszystkich} domkniętych podzbiorów \(X\),
włącznie ze zbiorem pustym.
Dla otwartego \(U\subseteq X\) i zwartego
\(C\subseteq X\) definiujemy
\[
 U^-_{\!F}=\{F\in\mathcal F(X):F\cap U\neq\varnothing\},
 \qquad
 C^0=\{F\in\mathcal F(X):F\cap C=\varnothing\}.
\]
\emph{Topologia Fella} jest topologią generowaną
przez wszystkie takie zbiory.
Pierwszy rodzaj oznacza trafienie otwartego zbioru,
drugi zaś ominięcie zbioru zwartego.
\end{definicja}

\begin{uwaga}[Konwencja dotycząca pustego zbioru]
\label{rem:topog21-fell-empty}
Niektóre źródła rozważają wyłącznie niepuste
zbiory domknięte. W tym rozdziale
\(\mathcal F(X)\) zawiera \(\varnothing\);
to założenie jest istotne dla
twierdzenia o zwartości poniżej.
Przy \(X\) niezwartym ciąg niepustych
zbiorów może w topologii Fella zbiegać
do \(\varnothing\).
\end{uwaga}

\begin{lemat}[Kryterium zwartości przez podbazę]
\label{lem:topog21-alexander}
Niech \(Y\) będzie przestrzenią topologiczną,
a \(\mathcal S\) podbazą jej topologii.
Jeżeli każde pokrycie \(Y\) zbiorami z \(\mathcal S\)
ma skończone podpokrycie, to \(Y\) jest zwarta.
\end{lemat}
\begin{proof}
Niech \(\mathcal U\) będzie dowolnym ultrafiltrem na \(Y\).
Gdyby nie był zbieżny do żadnego punktu,
dla każdego \(y\in Y\) istniałoby otwarte
otoczenie \(O_y\ni y\), które nie należy do \(\mathcal U\).
Ponieważ skończone przecięcia elementów podbazy
tworzą bazę, możemy wybrać
\(y\in S_{y,1}\cap\cdots\cap S_{y,m_y}\subseteq O_y\),
gdzie \(S_{y,i}\in\mathcal S\).
Co najmniej jeden z elementów \(S_{y,i}\)
nie należy do ultrafiltru: gdyby wszystkie
należały, zawierałby też ich przecięcie
i jego nadzbiór \(O_y\).
Wybierzmy taki element i oznaczmy go \(S_y\).
Rodzina \((S_y)_{y\in Y}\) pokrywa \(Y\).
Z założenia mamy skończone podpokrycie
\(S_{y_1},\ldots,S_{y_m}\).
Ultrafiltr zawiera co najmniej jeden
element skończonego pokrycia całej przestrzeni,
sprzecznie z wyborem \(S_y\).
Każdy ultrafiltr jest więc zbieżny,
a Twierdzenie~\ref{thm:topog-zwartosc-ultrafiltry}
kończy dowód.
\end{proof}

\begin{twierdzenie}[Zwartość topologii Fella]
\label{thm:topog21-fell-compact}
Dla dowolnej przestrzeni topologicznej \(X\)
przestrzeń \(\mathcal F(X)\) z topologią Fella
jest zwarta.
\end{twierdzenie}
\begin{proof}
Na mocy Lematu~\ref{lem:topog21-alexander}
wystarczy rozważyć pokrycie podbazowe
\[
 \mathcal F(X)=
  \bigcup_{i\in I}(U_i)^-_{\!F}\ \cup\
  \bigcup_{j\in J}C_j^0,
\]
gdzie \(U_i\) są otwarte, a \(C_j\) zwarte,
i pokazać, że ma skończone podpokrycie.
Przypuśćmy przeciwnie, że takiego podpokrycia
nie ma. Połóżmy
\(F=X\setminus\bigcup_{i\in I}U_i\).
Jest to zbiór domknięty, także gdy \(F=\varnothing\),
więc \(F\in\mathcal F(X)\), i żaden
\((U_i)^-_{\!F}\) go nie zawiera.

Dla ustalonego \(j\in J\) i każdego skończonego
\(I_0\subseteq I\) rodzina
\(\{(U_i)^-_{\!F}:i\in I_0\}\cup\{C_j^0\}\)
nie pokrywa \(\mathcal F(X)\).
Istnieje więc domknięty \(E\subseteq X\),
który nie trafia żadnego \(U_i\), \(i\in I_0\),
a jednocześnie trafia \(C_j\).
To znaczy
\[
 C_j\cap\bigcap_{i\in I_0}(X\setminus U_i)
 \neq\varnothing.
\]
Rodzina zbiorów \(C_j\cap(X\setminus U_i)\),
\(i\in I\), ma własność skończonych przecięć.
Ponieważ \(C_j\) jest zwarty, mamy
\(C_j\cap F\neq\varnothing\).
Zatem \(F\notin C_j^0\) dla każdego \(j\in J\).
Zbudowaliśmy element \(\mathcal F(X)\)
nienależący do żadnego zbioru pokrycia,
sprzeczność. Każde pokrycie podbazowe
ma zatem skończone podpokrycie.
\end{proof}

\begin{twierdzenie}[Oddzielanie w topologii Fella]
\label{thm:topog21-fell-hausdorff}
Jeżeli \(X\) jest lokalnie zwartą przestrzenią
Hausdorffa, to \(\mathcal F(X)\) z topologią Fella
jest zwartą przestrzenią Hausdorffa.
\end{twierdzenie}
\begin{proof}
Zwartość wynika z
Twierdzenia~\ref{thm:topog21-fell-compact}.
Niech \(F,G\in\mathcal F(X)\) będą różne.
Po zamianie ról możemy przyjąć, że
istnieje \(x\in F\setminus G\).
Zbiór \(X\setminus G\) jest otwartym
otoczeniem \(x\). Na mocy
Lematu~\ref{lem:topog-baire-zwarte-domkniecie}
istnieje otwarty \(W\ni x\) o zwartym
domknięciu \(C=\overline W^{\,X}\subseteq X\setminus G\).
Wtedy \(W^-_{\!F}\) jest otwartym otoczeniem \(F\),
a \(C^0\) otwartym otoczeniem \(G\).
Są one rozłączne, ponieważ \(W\subseteq C\).
\end{proof}

\begin{stwierdzenie}[Zbieżność Fella przez trafianie i omijanie]
\label{prop:topog21-fell-nets}
Niech \(X\) będzie przestrzenią topologiczną,
\((F_\lambda)_{\lambda\in D}\) siecią
w \(\mathcal F(X)\), a \(F\in\mathcal F(X)\).
Sieć \(F_\lambda\to F\) w topologii Fella
dokładnie wtedy, gdy spełnione są dwa warunki:
\begin{enumerate}[label=\textup{(\roman*)}]
\item dla każdego otwartego \(U\subseteq X\)
z \(F\cap U\neq\varnothing\)
od pewnego indeksu \(F_\lambda\cap U\neq\varnothing\);
\item dla każdego zwartego \(C\subseteq X\)
z \(F\cap C=\varnothing\)
od pewnego indeksu \(F_\lambda\cap C=\varnothing\).
\end{enumerate}
\end{stwierdzenie}
\begin{proof}
Są to dokładnie warunki, by sieć ostatecznie
należała do każdego podbazowego otoczenia \(F\).
Skończone przecięcie takich otoczeń
również jest ostatecznie trafiane:
dla skończonej liczby indeksów progowych
wybieramy wspólne górne ograniczenie
w skierowanym zbiorze \(D\).
Otrzymujemy zbieżność w topologii
generowanej przez podbazę.
\end{proof}

\section{Porównanie topologii Vietorisa i Fella}
\label{sec:topog21-comparison}

\begin{twierdzenie}[Zgodność w zwartej przestrzeni Hausdorffa]
\label{thm:topog21-fell-vietoris}
Jeżeli \(X\) jest zwartą przestrzenią Hausdorffa,
to na \(\mathcal K(X)=\mathcal F(X)\setminus\{\varnothing\}\)
topologia Vietorisa pokrywa się z topologią
indukowaną przez topologię Fella.
\end{twierdzenie}
\begin{proof}
W przestrzeni zwartej Hausdorffa
każdy niepusty zbiór domknięty jest zwarty,
więc obie rozważane rodziny punktów są równe.
Warunki dolne są identyczne.
Dla otwartego \(U\subseteq X\)
jego dopełnienie \(X\setminus U\) jest zwarte, zatem
\[
 U^+=(X\setminus U)^0
          \cap\mathcal K(X).
\]
Odwrotnie, dla dowolnego zwartego \(C\subseteq X\)
własność Hausdorffa zapewnia domkniętość \(C\),
a więc \(X\setminus C\) jest otwarty i
\[
 C^0\cap\mathcal K(X)=(X\setminus C)^+.
\]
Obie podbazy indukują tę samą topologię.
\end{proof}

\begin{przyklad}[Punkty uciekające do nieskończoności]
\label{ex:topog21-escaping}
Niech \(X=\mathbb R\).
W \(\mathcal F(\mathbb R)\) z topologią Fella
zachodzi \(\{n\}\to\varnothing\).
Istotnie, zbiór pusty nie ma otoczeń wymagających
trafienia jakiegokolwiek otwartego \(U\);
dla każdego zwartego \(C\subseteq\mathbb R\)
od pewnego \(n\) mamy \(n\notin C\).
Nie ma natomiast granicy tego ciągu
w \(\mathcal K(\mathbb R)\) z topologią Vietorisa,
bo \(d_H(\{n\},K)\to\infty\)
dla każdego ustalonego \(K\in\mathcal K(\mathbb R)\).

Ponadto \(K_n=\{0,n\}\) zbiega w topologii Fella
do \(K=\{0\}\): każdy otwarty zbiór trafiany
przez \(K\) jest trafiany przez \(K_n\),
a każdy zwarty \(C\) rozłączny z \(K\)
jest ostatecznie omijany przez \(K_n\),
ponieważ \(n\) opuszcza \(C\).
Jednak
\[
 d_H(K_n,K)=n
\]
dla \(n\ge1\), więc \(K_n\) nie zbiega
do \(K\) w topologii Vietorisa.
W szczególności na \(\mathcal K(\mathbb R)\)
topologia Fella indukowana z \(\mathcal F(\mathbb R)\)
jest ściśle słabsza od topologii Vietorisa.
\end{przyklad}

\begin{uwaga}[Dokładna rola zwartości]
\label{rem:topog21-restriction}
Dla dowolnej przestrzeni Hausdorffa \(X\)
ograniczenie podbazy Fella do \(\mathcal K(X)\)
jest zawarte w topologii Vietorisa:
dla zwartego \(C\) zbiór \(C\) jest domknięty,
więc warunek ominięcia \(C\) jest warunkiem
zawierania się w otwartym \(X\setminus C\).
Odwrotne zawieranie może zawieść,
jak pokazuje poprzedni przykład.
Nie wolno utożsamiać topologii Fella
z metryką Hausdorffa na niezwartym \(X\).
\end{uwaga}

\section{Podsumowanie}
\label{sec:topog21-summary}

W topologii Vietorisa punktami są niepuste
zbiory zwarte; wymaga się trafiania wskazanych
otwartych podzbiorów oraz zawierania w otwartym
otoczeniu całego zbioru. Dla przestrzeni metrycznej
topologia ta jest wyznaczana przez metrykę Hausdorffa.
Zupełność \(X\) przechodzi na \(\mathcal K(X)\);
zwartość \(X\) pociąga zwartość \(\mathcal K(X)\)
zarówno w przypadku metrycznym,
jak i w ogólnym przypadku zwartym Hausdorffa.

W topologii Fella punktami są wszystkie
zbiory domknięte, także pusty; wymagamy
trafiania otwartych i omijania zwartych
podzbiorów. Powstała przestrzeń jest zwarta
dla dowolnego \(X\), a Hausdorffa,
gdy \(X\) jest lokalnie zwarta i Hausdorffa.
Dla zwartej przestrzeni Hausdorffa
topologie Fella i Vietorisa zgadzają się
na niepustych zbiorach domkniętych.
Dla niezwartych przestrzeni zbiór może
utracić odległe punkty w granicy Fella,
co nie zachodzi przy zbieżności
w metryce Hausdorffa.